\begingroup
\begin{center}\textbf{Resumen}\end{center}
\fontsize{9.1}{10.55}\selectfont
\setlength{\parskip}{1.2pt}
\setlength{\abovedisplayskip}{3pt}
\setlength{\belowdisplayskip}{3pt}
Este artículo demuestra afirmativamente la hipótesis del continuo en la
Holografía Modular Triádica (HMT). El dominio no se inicia en una recta real
ya dada. Se construye mediante la composición de All Possible Paths (APP), la
firma ternaria orientada
\(\mathrm{TRIT}=\{-1,0,+1\}\) y el Topological Phase Kernel (TPK), o núcleo
topológico de fase. Su límite inverso
\(X_\infty^{\mathrm{enr}}\) conserva valor, hoja, orientación, residuo,
cociente, acarreo, frontera, ruta, memoria y supervivencia. La estructura
discreta conjunta del continuo y su recta arquimediana son salidas de esa
acción, no datos empleados para seleccionar estados.

Sea \(h\) el código efectivo de los grafos de los operadores APP, TRIT y TPK,
de sus prolongaciones, de sus restricciones y de las dos proyecciones. Para
cada historia realizada \(m_\infty\), la clausura genealógica
\(\mathfrak G_{\mathrm{HMT}}(h,m_\infty)\) se obtiene mediante iteración
transfinita de definiciones finitas con parámetros de niveles anteriores. El
teorema principal establece
\[
 \boxed{\mathfrak G_{\mathrm{HMT}}(h,m_\infty)
 \models 2^{\aleph_0}=\aleph_1}
\]
y, para todo
\(A\in\mathcal P^{\mathfrak G}
(\mathbb R_{\mathrm{HMT}}^{\mathfrak G})\), prueba la dicotomía
\[
 |A|\leq\aleph_0
 \qquad\text{o}\qquad
 |A|=|\mathbb R_{\mathrm{HMT}}^{\mathfrak G}|.
\]
Por tanto, no existe un cardinal intermedio entre los naturales y la recta
dentro del continuo íntegro generado por HMT.

La demostración encadena cuatro pasos. La cobertura exhaustiva de cilindros
identifica la recta publicada con
\(\mathcal P^{\mathfrak G}(\omega)\). La condensación genealógica construye
una inyección
\(j_{\mathrm{HMT}}:\mathbb R_{\mathrm{HMT}}^{\mathfrak G}
\hookrightarrow\omega_1^{\mathfrak G}\). La codificación de buenos órdenes
numerables y su elevación ternaria construyen la inyección en sentido contrario
\(k_{\mathrm{HMT}}\). Cantor--Bernstein concluye la igualdad cardinal, y el
argumento de cofinalidad extiende la conclusión a todos los subconjuntos de la
clausura completa. Ninguno de estos pasos presupone la igualdad que demuestra.

La fibra local de \(3^6=729\) subcilindros de refinamiento no acota el
resultado: la ley coinductiva actúa a toda profundidad prescrita. El retorno
nonádico recupera la fase mientras avanzan acarreo y memoria; no reinicia el
estado. La misma genealogía produce, como publicaciones correlacionadas, la
holonomía \(\Gamma_9\), el registro dodecafásico integral \(K\), la
cuatrirrelación \((\pi,\varphi,e,\alpha)\) y la realización incidencial
Paley--Witt--Golay--Leech--\(V^\natural\)--Monstruo. Esas salidas no prueban
la igualdad cardinal por coincidencia numérica; certifican la unidad del
estado cuya doble proyección publica valor y conserva estructura.

Finalmente,
\(\mathfrak G_{\mathrm{HMT}}(h,m_\infty)=L[h,m_\infty]\) se demuestra como
reconocimiento posterior de la jerarquía ya construida. No genera la recta,
no selecciona las inyecciones y no interviene en la reducción de cilindros.

\medskip
\noindent\textbf{Palabras clave:} Holografía Modular Triádica; APP; TRIT;
TPK; estado enriquecido; estructura discreta del continuo; clausura
genealógica; hipótesis del continuo; cardinalidad; ordinalidad; coinducción.
\endgroup
